\subsubsection{ISP, হোস্টিং এবং ফাইলের চলার দিক}
\textbf{ISP (Internet Service Provider)} ব্যবহারকারীকে ইন্টারনেট সংযোগ দেয়। \textbf{হোস্টিং সেবা} ওয়েবসাইটের ফাইল ও প্রয়োজনীয় সার্ভারের ব্যবস্থা দেয়। একই প্রতিষ্ঠান দুটো সেবা দিতে পারে, কিন্তু কাজ দুটি আলাদা। বাড়ির ইন্টারনেট সংযোগ থাকলেই কলেজের ওয়েবসাইট প্রকাশিত হয়ে যায় না; প্রকাশযোগ্য ফাইল সার্ভারে দিতে হয়।

\par\noindent\begin{minipage}[t]{.48\linewidth}\vspace{0pt}\small
\textbf{\color{chFourTeal}আপলোড (Upload)}\par
নিজের যন্ত্র থেকে দূরের সার্ভারে ফাইল পাঠানো। শিক্ষক নিজের কম্পিউটার থেকে \texttt{routine.pdf} সাইটে পাঠালে সেটি আপলোড।
\end{minipage}\hfill
\begin{minipage}[t]{.48\linewidth}\vspace{0pt}\small
\textbf{\color{chFourTeal}ডাউনলোড (Download)}\par
দূরের সার্ভার থেকে নিজের যন্ত্রে ফাইল নেওয়া। শিক্ষার্থী সেই \texttt{routine.pdf} ফোনে সংরক্ষণ করলে সেটি ডাউনলোড।
\end{minipage}\par
দিকটি ব্যবহারকারীর যন্ত্রের দৃষ্টিতে বোঝানো হয়েছে। ফাইল আদান-প্রদান শুধু FTP দিয়ে হয় না; HTTPS-এও ওয়েবসাইটে ফাইল পাঠানো ও নেওয়া যায়।

\subsubsection{পূর্ণ ডোমেইন নাম ও নাম-ঠিকানার সমন্বয়}
DNS-এর মূল স্তর পর্যন্ত সব লেবেলসহ সম্পূর্ণ নামকে FQDN (Fully Qualified Domain Name) বলে। যেমন, \texttt{www.example.com.}; শেষের ডটটি মূল স্তর বোঝায়। সাধারণ ব্যবহারে সেই ডট বাদ দিয়ে \texttt{www.example.com} লেখা হয়। FQDN কোনো নির্দিষ্ট পেজের পূর্ণ URL নয়; এতে \texttt{https://} বা ফাইলের পথ থাকে না।

\textbf{IANA (Internet Assigned Numbers Authority)} ইন্টারনেটের আইপি ঠিকানার পরিসর, DNS-এর মূল স্তর ও প্রোটোকলের নির্দিষ্ট সংখ্যার বৈশ্বিক সমন্বয় করে। আইপি ঠিকানার বড় পরিসর আঞ্চলিক নিবন্ধন সংস্থায় বরাদ্দ হয়; পরে সংশ্লিষ্ট সংস্থা ও সেবাদাতার মাধ্যমে ব্যবহারকারীর কাছে পৌঁছায়। তাই IANA প্রতিটি ব্যবহারকারীর যন্ত্রে নিজে ঠিকানা বসিয়ে দেয় না।

\textbf{ICANN (Internet Corporation for Assigned Names and Numbers)} ইন্টারনেটের স্বতন্ত্র নাম ও সংখ্যাভিত্তিক পরিচয়ের সমন্বয়ে কাজ করে। সাধারণ ব্যবহারকারী নিবন্ধন সেবাদাতার মাধ্যমে ডোমেইন নিবন্ধন করেন। ডোমেইন নিবন্ধন, আইপি ঠিকানার বরাদ্দ ও ওয়েব হোস্টিং সম্পর্কিত হলেও এক কাজ নয়।
